% Propietario conservado: /Users/ruben/Documents/New project/output/REVISION_ARGUMENTAL_APERTURAS_CIERRES_20260915/IX/ES/sections/ch_profundidad_constantes.tex
\section{Coinductive depth of the fourfold relation}
\label{ix:sec:profundidad-cuatrirrelacion-ch}

The regional readers and dodecaphase closure publish, respectively,
\(\pi,\varphi,e\) and \(\alpha\) from the common enriched state. We reserve
\(N\) for operator depth and \(n\) for decimal length:
the two indices are not identified and no uniform reading cost is assumed.
For each exact realization \(\chi\), its integer part is preserved
separately and \(p_{\chi,n}\) denotes the prefix of length \(n\) of its
fractional part, in the canonical decimal expansion not eventually equal
to nine. Define
\[
 \rho^{\rm dec}_{n+1,n}(d_1\ldots d_nd_{n+1})=d_1\ldots d_n,
 \qquad p_{\chi,0}:=\varnothing.
\]
The uniqueness of \(p_{\pi,n},p_{\varphi,n},p_{e,n}\) follows from the
uniqueness of the exact character already generated and of its canonical
decimal cylinder. That of \(p_{\alpha,n}\) follows from the integral
dodecaphase closure, its tail identity, and the projective compatibility
proved in Theorem~\ref{ix:thm:alpha-salida-dodecafasica-completa}. The
ninth-order analytic resolvent provides an additional finite
characterization, while the complete analytic reader computes the remaining
coefficients by \(b_{10+3n}=\lambda/729^n\), proves uniform convergence, and
constructs rational intervals contained in the integral cylinders. It thereby
identifies that same section at arbitrary depth by
Theorem~\ref{ix:prop:alpha-dos-limites}. The finite jet is used neither as a
generator nor as a substitute for the coinductive prolongation.

\begin{theorem}[Projective compatibility of \(\pi,\varphi,e,\alpha\) at arbitrary depth]
\label{ix:thm:profundidad-cuatrirrelacion-ch}
The coinductive action of the common state determines
\[
 \begin{gathered}
  \forall\chi\in\{\pi,\varphi,e,\alpha\}\;
  \forall n\ge1\quad\exists!\,p_{\chi,n},\\
  \rho^{\rm dec}_{n+1,n}(p_{\chi,n+1})=p_{\chi,n},
  \qquad
  p_{\chi,\infty}:=(p_{\chi,n})_{n\ge0}
  \in\varprojlim_{n\ge0}\{0,\ldots,9\}^{n}.
 \end{gathered}
\tag{D1}
\label{ix:eq:profundidad-cuatrirrelacion-ch}
\]
In particular,
\[
 (p_{\pi,\infty},p_{\varphi,\infty},p_{e,\infty},p_{\alpha,\infty})
\tag{D2}
\label{ix:eq:cuatrirrelacion-infinita-ch}
\]
is a correlated coinductive publication of the unique enriched state.
\end{theorem}

\begin{proof}
Result~\ref{ix:gen:prefijos-regionales} fixes the exact realizations
of the three regional readers.
Theorem~\ref{ix:thm:alpha-salida-dodecafasica-completa} constructs the complete
section \(\alpha_{\mathrm{HMT}}\) and proves compatibility of all its
prefixes. The canonical decimal expansion of each of these
four realizations has a unique prefix of each length; truncating the
prefix of length \(n+1\) gives the one of length \(n\). Applying those readers
to the common state gives the stated correlated family. Existence and
compatibility of this exact publication do not identify \(n\) with \(N\)
or add a cost bound to the reading procedures already established.
\end{proof}

The internal constructive order preserves the readers of
\(\pi,\varphi,e\), the signed state, the integer dodecaphase seal \(K\) and the
dodecaphase wheel that publishes \(\alpha\); no conventional constant is
promoted to a generator. An execution at one thousand, ten thousand or one million digits
verifies the corresponding finite computational instance of
\eqref{ix:eq:profundidad-cuatrirrelacion-ch}; it does not constitute a bound on the
theorem. The theorem is the compatible family for every \(n\) and the infinite
section determined by it.

This depth is not used to infer the continuum hypothesis. The four publications
certify the coinductive action of the same generator; cardinal closure has already been
obtained through coverage, condensation and the two injections. Its
verification criterion is twofold: compatibility of all truncations
of \eqref{ix:eq:profundidad-cuatrirrelacion-ch} and internal provenance of the
data of the dodecaphase wheel that publishes \(\alpha_{\rm HMT}\).
